\subsection{Philosophical Foundation: Event-First as Ontological Commitment}\label{subsec:philosophical-foundation}

ADL Lite inverts the object-first stance of conventional operational ontologies \cite{palantir_fde,dendron2020}, following the programme of formal ontology as an analytical tool for information systems \cite{guarino_1998}. Drawing on Wittgenstein's \emph{Tractatus} \S1.1---\emph{``Die Welt ist die Gesamtheit der Tatsachen, nicht der Dinge''} \cite{wittgenstein_tractatus,event_sourcing}---the system treats the occurrent as a fundamental primitive, and treats lifecycle statuses as institutional facts whose reality is constituted by recorded events \cite{searle1995construction}. This yields three governing principles: (1)~\textbf{No implicit state changes.} Every mutation is an explicit event. (2)~\textbf{Events are occurrents.} They happen, but do not endure \cite{bfo,dolce}. (3)~\textbf{All derived properties are historical.} Status, confidence, and validators are computed by deterministic functions over the event sequence and are never stored.

\subsection{Categories and Taxonomy}\label{subsec:categories}

We analyze ADL Lite's categories with reference to BFO \cite{bfo}, DOLCE \cite{dolce}, and UFO \cite{ufo}. Table~\ref{tab:upper-ontology-mapping} presents the cross-mapping.

\begin{table}[htbp]
\centering
\small
\caption{Cross-mapping of ADL Lite categories to upper ontologies. Entries marked ``---'' indicate intentional non-alignment; those marked ``\texorpdfstring{$\approx$}{~}'' denote approximate correspondence.}
\label{tab:upper-ontology-mapping}
\begin{tabularx}{\textwidth}{@{}p{1.4cm}p{2.5cm}XXXX@{}}
\toprule
\textbf{ADL Lite} & \textbf{BFO 2.0} & \textbf{DOLCE} & \textbf{UFO} & \textbf{Note} \\
\midrule
Event & \texttt{occurrent} & \texttt{perdurant} & \texttt{Event} & Temporal entity; happens; no endurance \\
EventChain & \texttt{process} & \texttt{accomplishment}$\approx$ & \texttt{Complex Event} & Ordered event sequence; \emph{process} layer (see record layer below) \\
EventChain-record & \texttt{information content entity} & \texttt{information object} & \texttt{Information Entity} & Serialized ICE; \emph{two-level account} in \S\ref{subsec:two-level-account} \\
Concept & \texttt{generically dependent continuant} & \texttt{non-physical object}$\approx$ & \texttt{Conceptual Entity} & Depends on EventChain-record as ICE bearer (not occurrent) \\
Relation (L3) & \texttt{relational quality} & \texttt{quality} & \texttt{Relator} & Mediates between concepts \\
Action (L4) & \texttt{planned process} & \texttt{action} & \texttt{Action} & Intentional occurrent with preconditions \\
Actor & \texttt{agent role} & \texttt{physical agent}$\approx$ & \texttt{Agent} & Role realized in material entity (human or software)\footnote{In BFO, agency is a role that inheres in a material entity (human or software process) and is realized in processes (events). We map ADL Lite Actor to BFO \texttt{agent role} because LLM agents are software processes that realize the agent role through event participation. This avoids the category error of treating software as a physical agent.} \\
Payload & \texttt{information content entity} & \texttt{information object} & \texttt{Information Entity} & Structured data carried by event \\
Hash & \texttt{generically dependent continuant} & --- & --- & Cryptographic identifier; no physical bearer \\
\bottomrule
\end{tabularx}
\end{table}

\subsubsection{Event as Occurrent/Perdurant}

In BFO, an \emph{occurrent} unfolds in time and is ontologically distinct from \emph{continuants} \cite{bfo}; in DOLCE, the corresponding category is \emph{perdurant} \cite{dolce}. ADL Lite's \texttt{Event} is both: a temporal entity with a definite beginning and end, no spatial location, and no endurance. The serialized record is an information content entity (ICE), not the event itself. In UFO, \texttt{REGISTER} and \texttt{VALIDATE} are \emph{actions} (intentional), whereas \texttt{RELATE}, \texttt{EVIDENCE}, and \texttt{SEAL} are \emph{events} (structural changes) \cite{ufo}.

\subsubsection{Concept as Generically Dependent Continuant}

An ADL Lite concept is a \emph{generically dependent continuant} (GDC) in BFO terms \cite{bfo}: it depends on some continuant for its existence, but not on any \emph{specific} one. Following the standard BFO/IAO analysis of information artifacts \cite{iao,smith2012ontology}, the dependence chain is:

\begin{center}
\textsf{Concept} (GDC) $\xrightarrow{\text{g-depends on}}$ \textsf{EventChain-record} (ICE) $\xrightarrow{\text{concretized\_by}}$ \textsf{Markdown\_file} (IBE)
\end{center}

This is the canonical BFO pattern: a GDC (the concept) generically depends on an information content entity (ICE, the serialized chain record), which is itself concretized by an information-bearing entity (IBE, the Markdown file or database row) \cite{smith2012ontology,limbaugh2024cyber}. The GDC does not depend on the occurrent process of events (the EventChain-process), because BFO's axiom D5 prohibits GDCs from depending on occurrents \cite{bfo}. The concept's causal origin is the genesis \texttt{REGISTER} event (an occurrent), but its ontological bearer is the ICE record (a continuant)---a distinction between \emph{causal origin} and \emph{ontological dependence} that Fine \cite{fine1995ontologicaldependence} and BFO both endorse. The concept's identity is determined by the \emph{genesis event} (the first event in the chain), not by the totality of events---analogous to a poem, which depends on some physical copy for its existence but not on any specific one. In DOLCE, the closest category is \emph{non-physical object} \cite{dolce}; in UFO, concepts are \emph{conceptual entities} \cite{ufo}.

\subsubsection{EventChain as Process}

An EventChain is a \emph{process} in BFO: a complex occurrent with temporal parts (the individual events) that unfolds over time \cite{bfo}. It is not a continuant because it does not exist in full at any moment. It is individuated by the ordered sequence of events and their cryptographic linkage (see Section~\ref{subsec:two-level-account} for the full two-level account, which distinguishes the process from its record).

\subsubsection{Actor as Agent}

ADL Lite's \texttt{Actor} is mapped to BFO \texttt{agent role} and UFO \texttt{Agent}; the DOLCE \texttt{physical agent} mapping is approximate because DOLCE lacks a non-physical agent category (see Table~\ref{tab:upper-ontology-mapping} footnote for full rationale). Future work (FW1) may define a DOLCE extension for software agents.

\subsubsection{Relation as Relational Quality / Relator}\label{subsec:relation-relator}

ADL Lite's L3 relation blocks assert typed edges between concepts. In BFO, these are \emph{relational qualities}: qualities that inhere in one entity and bear a relation to another \cite{bfo}. ADL Lite's L3 relation block is inscribed in the source concept's Markdown file (the bearer), while the relation mediates between two concepts. This reading is consistent with BFO relational qualities, which need not be symmetric. In UFO, the same relations are \emph{relators}: entities that mediate between relata and are brought into existence and destroyed by events \cite{ufo}. The dual reading---BFO relational quality for the record, UFO relator for the mediation---reflects the two-level account of Section~\ref{subsec:two-level-account}.

\textbf{Creation.} A \texttt{RELATE} event $e$ establishes a relator $r$ at time $t = e.\text{timestamp}$:
\begin{multline}\label{eq:relation-creation}
    \text{HoldsAt}\bigl(\text{Relation}(c_1, c_2, p), t\bigr) \iff \exists e \in C.\; e.\text{type} = \text{RELATE} \land e.\text{timestamp} \leq t \\
    \land \neg\text{Revoked}(c_1, c_2, p, t)
\end{multline}
where $\text{Revoked}(c_1, c_2, p, t)$ holds if and only if there exists a revocation event $e'$ with $e'.\text{timestamp} \leq t$.

\textbf{Revocation.} A relator is terminated by an explicit \texttt{REVOKE} event $e_{\text{rev}}$:
\begin{multline}\label{eq:relation-revocation}
    \text{Revoked}(c_1, c_2, p, t) \iff \exists e_{\text{rev}} \in C.\; e_{\text{rev}}.\text{type} = \text{REVOKE} \land e_{\text{rev}}.\text{timestamp} \leq t \land \\
    e_{\text{rev}}.\text{payload}.\text{target\_concept\_id} = c_2 \land e_{\text{rev}}.\text{payload}.\text{predicate} = p
\end{multline}

\textbf{Note on revocation semantics.} ADL Lite supports two semantics for relation termination. \emph{Cessation semantics} (Equation~\ref{eq:relation-revocation}): a dedicated \texttt{REVOKE} event explicitly terminates the relator, whose record (L3 block + \texttt{REVOKE} event) remains as an auditable ICE. \emph{Epistemic weakening} (alternative): a \texttt{RELATE} event with \texttt{confidence}=0 suspends the relation's mediation without formally terminating it, allowing graded weakening. Importantly, \emph{HoldsAt does not depend on confidence}: it checks only whether a \texttt{REVOKE} event exists (Equation~\ref{eq:relation-revocation}). Confidence is epistemic strength, not an existence condition---a relation with confidence 0.3 still holds; a relation with a \texttt{REVOKE} event does not, regardless of the original \texttt{RELATE}'s confidence. We recommend \emph{cessation semantics} as the default for audit and compliance queries because it yields a binary, unambiguous existence condition; epistemic weakening is retained for consumer-side belief filtering (e.g., ranking by a confidence threshold $\theta$) and for systems preferring a graded thresholded variant $\text{HoldsAt}_\theta$. The \texttt{REVOKE} event is a communication event (not in $\Sigma_{\text{life}}$), so it does not affect $\delta(C)$; Theorems~1--6 are unchanged.

\textbf{Role constraints and temporal scoping.} A relation holds only while both relata exist and are neither deprecated nor archived (Equation~\ref{eq:relation-role-constraints}, where $\text{Deprecated}(c,t)$ and $\text{Archived}(c,t)$ are derived from $c$'s chain):
\begin{multline}\label{eq:relation-role-constraints}
    \text{HoldsAt}\bigl(\text{Relation}(c_1, c_2, p), t\bigr) \rightarrow \text{Exists}(c_1, t) \land \text{Exists}(c_2, t) \land \neg\text{Deprecated}(c_1, t) \land \\
    \neg\text{Deprecated}(c_2, t) \land \neg\text{Archived}(c_1, t) \land \neg\text{Archived}(c_2, t)
\end{multline}
Every relator has a temporal extent $[t_{\text{create}}, t_{\text{revoke}})$: it is a continuant that exists only during this interval; what endures beyond it is the record (the L3 block), an ICE documenting that the relation held. This exemplifies UFO's relator theory \cite{ufo}: existential dependence on both relata, event-based creation, event-based destruction, and mediation between entities rather than a property of either relatum alone; DOLCE's \emph{quality} category includes both intrinsic and relational qualities \cite{dolce}, and ADL Lite relations are relational qualities in this sense.

\subsubsection{Two-Level Ontological Account: Occurrents versus Records}\label{subsec:two-level-account}

The term ``EventChain'' is ambiguous between a process (events unfolding in time) and an ICE (the serialized record). These are distinct entities with distinct identity conditions.

\paragraph{Level 1: Occurrents (What Happens).}
Event ($e$): an individual occurrent/perdurant with a determinate timestamp. EventChain-process ($P$): the complex occurrent comprising the ordered event sequence. Action ($a$): an intentional occurrent. Level 1 entities do not endure; at time $t$, only events up to $t$ have occurred.

\paragraph{Level 2: Information Content Entities (What Is Recorded).}
EventChain-record ($R$): the serialized, hash-linked JSON/YAML representation---an ICE, a continuant existing in full when stored. Concept ($c$): the GDC depending on the record; it interprets the event sequence to project status, confidence, and validators. Relation ($r$): the record of a relational quality in L3 blocks, an ICE documenting a relationship that held during a specific interval.

\paragraph{Explicit identity axioms (necessary and sufficient conditions).}
Following Lowe's account of identity conditions as both necessary and sufficient criteria \cite{lowe1998possibility}, we state each axiom as a biconditional.

\begin{enumerate}[label=Axiom I\arabic*:,leftmargin=3em]
    \item \textbf{Event identity.} The identity of an event is \emph{necessary and sufficient} given by its UUID:
    \begin{equation}\label{eq:i1}
        e = e' \iff e.\text{event\_id} = e'.\text{event\_id}.
    \end{equation}
    The UUID is a rigid designator: if $e \neq e'$ then their event IDs differ, and identical event IDs entail the same event (no two distinct events may share a UUID).
    
    \item \textbf{Concept identity.} The identity of a concept is \emph{necessary and sufficient} given by its genesis hash:
    \begin{equation}\label{eq:i2}
        c = c' \iff c.\text{genesis\_hash} = c'.\text{genesis\_hash}.
    \end{equation}
    The \texttt{concept\_id} is a human-readable alias; two concepts with the same \texttt{concept\_id} but different genesis hashes are distinct registry items (see Section~\ref{subsec:identity-conditions} for registry-item versus domain-concept identity).
    
    \item \textbf{EventChain-record identity.} Two records are identical \emph{iff} they have the same concept and the same event sequence:
    \begin{equation}\label{eq:i3}
        R = R' \iff R.\text{concept\_id} = R'.\text{concept\_id} \land \text{events}(R) = \text{events}(R').
    \end{equation}
    Event-sequence equality is extensional: every event in $R$ appears in $R'$ with the same hash, and vice versa. The ordering is implied because the cryptographic hash chain enforces a unique total order.
    
    \item \textbf{EventChain-process identity.} Two processes are identical \emph{iff} they have the same events in the same order with the same timestamps:
    \begin{equation}\label{eq:i4}
        P = P' \iff \text{events}(P) = \text{events}(P') \land \text{order}(P) = \text{order}(P') \land \text{timestamps}(P) = \text{timestamps}(P').
    \end{equation}
    The timestamp condition is necessary because two processes with the same events in the same order but at different times are distinct occurrents (they occupy different temporal regions).
\end{enumerate}

\paragraph{FORK identity preservation.}
A \textsc{Fork} event creates a child concept $c'$ from a parent $c$. The child is a \emph{new} entity with a \emph{new} genesis hash:

\begin{enumerate}[label=Axiom I\arabic*:,leftmargin=3em,resume]
    \item \textbf{FORK identity.} Let $c'$ be the concept created by $\textsc{Fork}(c)$. Then:
    \begin{equation}\label{eq:i5}
        c' \neq c \;\land\; \text{genesis\_hash}(c') \neq \text{genesis\_hash}(c).
    \end{equation}
    The child's genesis hash is a function of the parent's genesis hash and the FORK event:
    \begin{equation}\label{eq:i5b}
        \text{genesis\_hash}(c') = \text{SHA-256}\bigl(\textsc{Fork\_event} \,\|\, \text{genesis\_hash}(c)\bigr).
    \end{equation}
    Even if the FORK event carries identical payload to the parent's genesis event, the child remains distinct because the FORK event has a unique \texttt{event\_id} and timestamp. This is \emph{provenance identity} (identity fixed by causal history) rather than content-addressed identity \cite{trusty_uris}.
\end{enumerate}

\paragraph{CRDT merge identity.}
When two branches $b_1, b_2$ of the same concept are merged, the result is a new record $R$ whose event set is the union of the two branch event sets:

\begin{enumerate}[label=Axiom I\arabic*:,leftmargin=3em,resume]
    \item \textbf{CRDT merge identity.} Let $R = \text{merge}(b_1, b_2)$. Then:
    \begin{equation}\label{eq:i6}
        \text{events}(R) = \text{sorted\_unique}\bigl(\text{events}(b_1) \cup \text{events}(b_2)\bigr),
    \end{equation}
    where sorting is deterministic by \texttt{event\_id} (UUID). The merge preserves the \texttt{concept\_id} of both parents:
    \begin{equation}\label{eq:i6b}
        R.\text{concept\_id} = b_1.\text{concept\_id} = b_2.\text{concept\_id}.
    \end{equation}
    Merge never produces a new \texttt{concept\_id}; it produces a new \emph{record} of the same concept. The identity of the concept $c$ is unchanged because the genesis hash (which fixes $c$'s identity, Axiom~I2) is preserved in both branches.
\end{enumerate}

\paragraph{L3 relation identity.}
An L3 relation is a relator that exists only during the interval $[t_{\text{create}}, t_{\text{revoke}})$. Its identity is fixed by the \textsc{Relate} event that created it:

\begin{enumerate}[label=Axiom I\arabic*:,leftmargin=3em,resume]
    \item \textbf{L3 relation identity.} Two relators are identical \emph{iff} they were created by the same \textsc{Relate} event:
    \begin{equation}\label{eq:i7}
        r = r' \iff \textsc{Relate\_event}(r).\text{event\_id} = \textsc{Relate\_event}(r').\text{event\_id}.
    \end{equation}
    A relator is terminated by a \textsc{Revoke} event (Equation~\ref{eq:relation-revocation}), after which it ceases to exist; the record of the relator (the L3 block) persists as an ICE, but the mediating entity no longer holds.
\end{enumerate}

\paragraph{Dependence axioms (necessary and sufficient conditions).}
\begin{enumerate}[label=Axiom D\arabic*:,leftmargin=3em]
    \item \textbf{Existential dependence on the record (D1).} A concept $c$ ontologically depends on the EventChain-record $R$: $c$ would not exist if $R$ had not been created. Formally, $c$ depends on $R$ iff $R$ is an EventChain-record about $c$:
    \begin{equation}\label{eq:d1}
        \text{Dep}(c, R) \iff \text{EventChainRecord}(R) \land \text{isAbout}(R, c).
    \end{equation}
    This is a generic dependence on an ICE bearer (D2), not a dependence on the EventChain-process.\footnote{We distinguish causal origin from ontological dependence. The \texttt{REGISTER} event is the \emph{causal origin} of $c$---it brought $c$ into existence---but $c$'s continued existence depends on the ICE record (D2), not on the occurrent process. This avoids the BFO constraint that GDCs cannot depend on occurrents (D5).}
    
    \item \textbf{Generic dependence (D2).} $c$ generically depends on the EventChain-record: it depends on an ICE bearer, but not on any specific physical copy. This is necessary but not sufficient for D1: generic dependence is the \emph{kind} of dependence, while D1 specifies the \emph{relatum}.
    
    \item \textbf{Concretization (D3).} The EventChain-record is concretized by a material bearer (file, database row, agent cache). This is specific dependence: the record depends on \emph{some} material bearer, and if that bearer is destroyed, the record ceases to exist unless another copy exists.
    
    \item \textbf{Process-to-record (D4).} The EventChain-process causally produces the EventChain-record. This is a causal, not ontological, dependence: the process is the \emph{origin} of the record, but the record's continued existence does not depend on the process continuing (the record may outlast the process).
    
    \item \textbf{No cross-level identity (D5).} No occurrent is identical to any ICE (a BFO constraint). Consequently, no GDC can ontologically depend on an occurrent; a GDC may only depend on continuants (including ICEs). Formally:
    \begin{equation}\label{eq:d5}
        \forall x \forall y\;\bigl(\text{GDC}(x) \land \text{Occurrent}(y) \rightarrow \neg\text{Dep}(x, y)\bigr).
    \end{equation}
    
    \item \textbf{Generic dependence of concept on record (D6).} Every concept generically depends on some EventChain-record that is about it:
    \begin{equation}\label{eq:d6}
        \forall c\;\bigl(\text{Concept}(c) \rightarrow \exists R\;(\text{EventChainRecord}(R) \land \text{isAbout}(R, c) \land \text{Dep}(c, R))\bigr).
    \end{equation}
    D6 is the universal closure of the left-to-right direction of D1: it states that \emph{every} concept must have \emph{some} record bearer. D6 is necessary for the existence of concepts; without a record, a concept cannot exist. D6 together with D2 (generic dependence) entails that a concept can survive the destruction of any particular record copy, provided at least one copy remains.
\end{enumerate}

\paragraph{First-order logic encoding.} To make the dependence structure amenable to formal analysis, we translate the axioms into first-order logic with explicit world parameters (following Fine's possible-worlds framework \cite{fine1995ontologicaldependence}), letting $E(x,w)$ mean ``$x$ exists in world $w$'' and $\text{Dep}(x,y,w)$ mean ``$x$ ontologically depends on $y$ in world $w$''. D2 (generic dependence) states that in every world where $c$ exists some record bears it, yet no \emph{particular} record is indispensable; D5 (no cross-level identity) states $\forall x \forall y\,(\text{Occurrent}(x) \land \text{ICE}(y) \rightarrow x \neq y)$, from which $\text{Concept}(c) \land \text{EventChainProcess}(P) \rightarrow \neg\text{Dep}(c,P,w)$ follows; D6 (universal record dependence) states that every concept depends on \emph{some record about it}; and I3--I4, I5, and I7 render record/process identity, FORK identity, and relator identity as extensional sorted-signature criteria. The complete FOL encodings (Equations~D2/D5/D6/I3--I5/I7) are given in Appendix~\ref{app:proofs}. The OWL~2 DL fragment (Appendix~\ref{app:owl-dl}) provides a decidable approximation of the class-level structure, while these FOL formulas capture the dependence and identity relations that exceed OWL expressivity.

\paragraph{Resolving the apparent ambiguity.}
$\text{EventChain-process}$ and $\text{EventChain-record}$ are distinct (Axioms I3--I4); the process causally produces the record (D4), which bears the concept (D2). A GDC cannot ontologically depend on an occurrent (D5), but can depend on the ICE that documents it (D2--D3). The concept's causal origin is traceable to the genesis event (the \texttt{REGISTER} occurrence), but its ontological bearer is the record. This resolves the category mistake.
\subsubsection{Action as Planned Process / Intentional Event}

ADL Lite's L4 action types are \emph{planned processes} in BFO (processes realizing a plan) \cite{bfo} and \emph{actions} in UFO (intentional events performed by agents to achieve goals) \cite{ufo}. The precondition system captures the plan aspect: an action can only be performed if certain conditions hold (e.g., \texttt{VALIDATE} requires status \texttt{provisional}), reflecting ontological norms governing which events are permissible in which states.

\subsection{OntoClean Evaluation}\label{subsec:ontoclean}

We evaluate ADL Lite's core classes through the OntoClean meta-property lens \cite{ontoclean,guarino2002evaluating,guarino2000identity}. OntoClean provides four meta-properties for ontological analysis: \emph{rigidity} ($R$), \emph{identity} ($I$), \emph{unity} ($U$), and \emph{dependence} ($D$) \cite{ontoclean}. A class is \emph{rigid} ($+R$) if its instances cannot cease to be members without ceasing to exist; \emph{carries identity} ($+I$) if its instances can be re-identified across worlds; \emph{carries unity} ($+U$) if its instances are unified wholes (not mereological sums); and \emph{existentially dependent} ($+D$) if its instances cannot exist without some other entity \cite{ontoclean,fine1995ontologicaldependence}. Supplementary Table S24 summarises the assignments for ADL Lite's ten core classes; the rigorous justification follows.

\subsubsection{Rigidity Analysis}\label{subsubsec:ontoclean-rigidity}

A class is rigid ($+R$) iff every instance is necessarily an instance in every possible world in which it exists \cite{ontoclean}. \texttt{Concept} is $+R$ because the genesis hash is an essential property (Axiom~I2): if a concept's genesis hash changed, it would be a different entity. \texttt{Hash} is $+R$ because a SHA-256 value is a fixed mathematical object. All other core classes are $-R$: an event may be re-interpreted without changing its UUID; the EventChain-process changes as events are appended; the record is replaced by a new snapshot with each append; actions are temporal slices that conclude; relations are created and terminated by \texttt{RELATE}/\texttt{REVOKE} events; payloads are transient; status values change as the chain evolves; and an actor is an agent role, which is anti-rigid by definition \cite{ontoclean}.

\subsubsection{Identity Analysis}\label{subsubsec:ontoclean-identity}

A class carries identity ($+I$) if instances can be re-identified across worlds \cite{ontoclean,guarino2000identity}. All ten classes carry identity ($+I$), because the event-first design requires every entity to be traceable: events by UUID (Axiom~I1); EventChain-processes by ordered events and timestamps (Axiom~I4); records by concept plus event set (Axiom~I3); concepts by genesis hash (Axiom~I2); relations by their creating \texttt{RELATE} event (Axiom~I7); actions by event UUID; actors by actor identifier; payloads by content hash; hashes by string equality; and status by concept-plus-time.

\subsubsection{Unity Analysis}\label{subsubsec:ontoclean-unity}

A class carries unity ($+U$) if its instances are unified wholes rather than mereological sums \cite{ontoclean}. \texttt{Concept}, \texttt{Relation}, \texttt{Actor}, and \texttt{EventChain}-record are $+U$: each is a coherent entity with a definite identity criterion (genesis hash, relator mediation, agenthood, structured schema). \texttt{Event}, \texttt{EventChain}-process, \texttt{Action}, \texttt{Payload}, \texttt{Status}, and \texttt{Hash} are $-U$: events are temporal slices, processes are mereological sums of events, and the rest are data structures or derived values.

\subsubsection{Dependence Analysis}\label{subsubsec:ontoclean-dependence}

A class is existentially dependent ($+D$) if its instances cannot exist without some other entity \cite{fine1995ontologicaldependence,ontoclean}. \texttt{Concept} is $+D$ (generic): as a GDC it depends on \emph{some} ICE bearer (BFO/IAO: GDC $\rightarrow$ ICE $\rightarrow$ IBE \cite{smith2012ontology,limbaugh2024cyber}), not on any specific copy. \texttt{EventChain}-record is $+D$ (generic, on a material bearer). \texttt{Relation} (L3) is $+D$ (specific, mutual rigid dependence on both relata). \texttt{Action}, \texttt{Payload}, and \texttt{Status} are $+D$ (specific: on actor, carrying event, and concept-plus-events respectively). \texttt{Event}, \texttt{EventChain}-process, \texttt{Actor}, and \texttt{Hash} are $-D$.

\subsubsection{Constraints and Consequences}\label{subsubsec:ontoclean-constraints}

The assignments satisfy OntoClean's logical constraints \cite{ontoclean,guarino2002evaluating}: $+R \Rightarrow +I$ (\texttt{Concept}, \texttt{Hash} both carry identity); $+U \Rightarrow +I$ (all $+U$ classes carry identity); and $+D \Rightarrow +I$ (all $+D$ classes carry identity). Subclass constraints also hold: \texttt{Concept} $\sqsubseteq$ \texttt{GDC} with both $+R$/$+U$/$+D$; \texttt{EventChain}-record $\sqsubseteq$ \texttt{ICE} with both $+U$/$+D$; \texttt{Relation} $\sqsubseteq$ \texttt{Relator} (UFO) with both $+U$/$+D$.

The assignments directly justify the design: (i)~the $+R$ assignment for \texttt{Concept} justifies genesis-hash identity; (ii)~the $+D$ assignments justify the BFO/IAO dependence chain; (iii)~forking creates a child with a new genesis hash, which is a new \emph{instance}, not a subclass, so rigidity is preserved; (iv)~the $-R$ assignment for \texttt{Event} justifies event immutability; and (v)~the $-R$ assignment for \texttt{Status} justifies status as a derived, non-essential property.

\subsubsection{Comparison with Alternative Readings}\label{subsubsec:ontoclean-alternatives}

Three alternative readings are inferior to the chosen assignments. \emph{Alternative 1 (Concept as independent continuant):} would make \texttt{Concept} $-D$, contradicting the BFO/IAO information-artifact analysis \cite{smith2012ontology} and importing spatiotemporal identity criteria inappropriate for abstract concepts. \emph{Alternative 2 (Event as continuant):} would let events endure and change, violating the event-first principle that all derived properties are historical (\S\ref{subsec:philosophical-foundation}) and undermining the append-only design. \emph{Alternative 3 (Status as rigid):} would make a \texttt{validated} and a \texttt{deprecated} concept different entities, contradicting explicit lifecycle transitions. These assignments reflect author judgment and would benefit from independent ontologist assessment (FW9).

\subsection{Identity Conditions}\label{subsec:identity-conditions}

We formalize identity conditions in Section~\ref{subsec:two-level-account} as \emph{necessary and sufficient} biconditionals: events by UUID (Axiom~I1, Equation~\ref{eq:i1}), concepts by genesis hash (Axiom~I2, Equation~\ref{eq:i2}), and chains by concept plus ordered events (Axioms~I3--I4, Equations~\ref{eq:i3}--\ref{eq:i4}). Here we examine the consequences for FORK, CRDT merge, and L3 relations. Deprecation and re-activation are status changes, not identity changes (the genesis hash is invariant). We evaluate the concept's identity through the OntoClean lens \cite{ontoclean}.

\textbf{Registry-item identity versus domain-concept identity.} We distinguish two senses of ``concept'' in ADL Lite. \emph{Registry-item identity} (operational) is fixed by the genesis hash; it is the identity criterion used by the EventChain data structure. Two registry entries with the same genesis hash are the same item, regardless of their content. \emph{Domain-concept identity} (intensional) is determined by the meaning, definition, and extensional content of the capability. Two registry items (different genesis hashes) may refer to the same domain concept (e.g., ``gradient-explosion'' and ``exploding-gradients'' in Section~\ref{subsec:domain-applicability}). The genesis-hash criterion provides \emph{registry-item} identity, not \emph{domain-concept} identity. It is intentionally operational: the registry must distinguish items even when their content is identical (to handle independent discovery by different agents). Domain-level identity is mediated by L3 relations (e.g., \texttt{isomorphic-to}, \texttt{specialisation-of}). This is analogous to a library catalog, in which each entry has a unique ISBN (registry identity) but multiple entries may refer to the same intellectual work (domain identity).

\textbf{Rigidity and sufficiency.} A concept's genesis hash is \emph{essential}---it is both necessary and sufficient for identity (Axiom~I2). It cannot change without the entity ceasing to exist, while its status is non-rigid. This aligns with OntoClean's recommendation to use rigid properties as identity criteria \cite{ontoclean}. The sufficiency direction ($c = c' \rightarrow \text{genesis\_hash}(c) = \text{genesis\_hash}(c')$) guarantees that distinct genesis hashes entail distinct concepts; the necessity direction guarantees that identical genesis hashes entail the same concept. This two-way criterion is stronger than a mere ``individuation'' principle; it is a full identity condition in Lowe's sense \cite{lowe1998possibility}.

\subsubsection{FORK Identity Preservation}\label{subsubsec:fork-identity}

A \textsc{Fork} event creates a child concept $c'$ from a parent $c$; the child is a \emph{new} entity with a \emph{new} genesis hash (Axiom~I5, Equation~\ref{eq:i5}), computed as $\text{genesis\_hash}(c') = \text{SHA-256}(\text{FORK\_event} \,\|\, \text{genesis\_hash}(c))$, where $\text{FORK\_event}$ includes a unique \texttt{event\_id}, timestamp, and \texttt{canon\_version}. Parent and child are \emph{never} identical, even for byte-identical payloads, because the FORK events have different UUIDs and timestamps. This embodies \emph{provenance identity} rather than content-addressed identity \cite{trusty_uris}: the FORK operation is a \emph{constructor} of new identity, and the \texttt{fork-of} L3 predicate records the lineage. If an agent modifies the payload during a fork, the child is still a new entity, with the modification recorded in the FORK payload; the parent's identity is unchanged. This satisfies the OntoClean constraint that rigid classes cannot be split into subclasses with different identity criteria.

\subsubsection{CRDT Merge Identity}\label{subsubsec:crdt-merge-identity}

When two branches $b_1, b_2$ of the same concept are merged, the result is a new record $R$ whose event set is the union of the two branch event sets, sorted deterministically by \texttt{event\_id} (Axiom~I6, Equation~\ref{eq:i6}); the merge preserves the \texttt{concept\_id} of both parents (Equation~\ref{eq:i6b}). Merge never produces a new \texttt{concept\_id} or genesis hash: it produces a new \emph{record} of the same concept (the identity of $c$ is unchanged because the genesis hash is preserved). Deterministic ordering by UUID (independent of any physical clock) is required for record identity: sorting by timestamp could produce different orderings under clock skew, violating $R = R' \iff \text{events}(R) = \text{events}(R')$ (Axiom~I3). Merge is defined only for branches of the same concept: attempts to merge branches with different \texttt{concept\_id}s are rejected by $\text{verify\_integrity}$ (Theorem~T1).

\subsubsection{L3 Relation Identity and Termination}\label{subsubsec:l3-relation-identity}

An L3 relation is a relator that exists only during $[t_{\text{create}}, t_{\text{revoke}})$; its identity is fixed by the \textsc{Relate} event that created it (Axiom~I7, Equation~\ref{eq:i7}). A relator $r$ between $c_1$ and $c_2$ \emph{mutually} and \emph{rigidly} depends on both concepts: if $c_1$ is archived, $r$ ceases to hold even if $c_2$ still exists. Formally, $\text{HoldsAt}(r, t) \rightarrow \text{Exists}(c_1, t) \land \text{Exists}(c_2, t) \land \neg\text{Archived}(c_1, t) \land \neg\text{Archived}(c_2, t)$ (a stronger condition than Equation~\ref{eq:relation-role-constraints}). A relator is terminated by a \textsc{Revoke} event (Equation~\ref{eq:relation-revocation}): after revocation it ceases to exist as a mediator, while the L3 block and the \textsc{Revoke} event persist as ICEs---an event-based destruction pattern consistent with UFO relator theory \cite{ufo}.

\subsubsection{Domain-Level Identity Alignment and Merge Conditions}\label{subsubsec:domain-identity-alignment}

While the genesis hash provides operational identity, domain-level identity (whether two chains denote the same concept) is determined by human curators, not the registry: ADL Lite does not support automatic chain merging (two chains with different genesis hashes are always distinct registry items), and domain-level consolidation is a governance decision recorded via L3 relations and \textsc{DEPRECATE} events. We recommend six normative L3 predicates, ordered from strongest to weakest claim: (i)~\texttt{isomorphic-to} (identical structure/content, different genesis hashes; reflexive, symmetric, transitive---the strongest claim short of genesis equality); (ii)~\texttt{specialisation-of} (a refinement or instance; irreflexive, antisymmetric, transitive); (iii)~\texttt{fork-of} (explicit lineage via a fork operation; irreflexive, antisymmetric, transitive by ancestry, stronger than \texttt{isomorphic-to} because it records provenance); (iv)~\texttt{related-to} (general association; symmetric, non-transitive); (v)~\texttt{analogical-to} (similar by analogy; symmetric, non-transitive); (vi)~\texttt{co-occurs-with} (frequent co-occurrence; symmetric, non-transitive).

\textbf{Relation semantics and merge conditions.} These properties are declared in \texttt{adl\_core\_ontology.yaml} but not enforced by the ActionExecutor at append time---enforcing transitivity would require a global graph search ($O(|V| + |E|)$ per \textsc{Relate} event), violating the $O(1)$ append complexity. Instead the query engine performs \emph{inference expansion} at read time using the transitive closure of \texttt{isomorphic-to} and \texttt{specialisation-of} (not \texttt{related-to} or \texttt{analogical-to}), separating \emph{storage semantics} (append-only, no global consistency check) from \emph{query semantics} (inference at read time). A domain expert may consolidate two chains by deprecating one in favour of the other via \texttt{isomorphic-to} or \texttt{specialisation-of}, recorded as a \textsc{DEPRECATE} event with reasoning---not an automatic merge---so the deprecated chain remains auditable and the full provenance of the consolidation is recoverable.

\textbf{Provenance-addressing vs. content-addressing.} The genesis hash is an \emph{operational (provenance) identifier}, not a content-addressed hash (Merkle DAG/IPFS CID): it includes the \texttt{event\_id} (UUID), \texttt{timestamp}, and \texttt{canon\_version} in addition to the payload, so two byte-identical \textsc{Register} events have different genesis hashes. This deliberately guarantees \emph{provenance uniqueness} (each registration is a distinct event with its own actor, timestamp, and UUID) rather than \emph{semantic equivalence}; content-addressing would enable automatic deduplication but lose the provenance needed to distinguish independent registrations. Domain-level consolidation of identical concepts is the responsibility of human curators via L3 predicates. Concept \emph{unity} is individuated by the genesis hash (the \texttt{EventChain} is a mereological sum of events with no unified whole beyond sequential aggregation \cite{ontoclean}), and the concept is \emph{generically dependent} on the EventChain record (an ICE bearer, not any specific physical copy \cite{bfo,fine1995ontologicaldependence}).

\subsection{Ontological Dependence}\label{subsec:ontological-dependence}

ADL Lite exhibits three forms of dependence. (1)~\emph{Identity and dependence conditions.} The concept is a GDC depending on the EventChain record: \textsf{Concept} $\xrightarrow{\text{GDC}}$ \textsf{ICE}(EventChain-record) $\xrightarrow{\text{concretized\_by}}$ \textsf{Markdown\_file}. A concept depends on the information artifact (the serialized record) rather than on the occurrent process of events, consistent with BFO's constraint that GDCs cannot depend on occurrents; its identity is fixed by the genesis event. As a GDC it can be concretized in multiple copies (Git repositories, local files, agent caches) without loss of identity, and it ceases to exist only when all concretizations are irreversibly destroyed \emph{and} no agent retains information to reconstruct the chain. (2)~\emph{Causal origin vs. ontological dependence.} A concept $c$ \emph{causally originates from} the EventChain-process (the genesis \textsc{Register} event brought it into existence) but \emph{ontologically depends} only on the EventChain-record $R$ (an ICE): the causal origin is an unrepeatable historical fact, whereas the ontological dependence is a repeatable bearer relationship (Axiom~D5 prohibits existential dependence on an occurrent). (3)~\emph{D6 (universal record dependence).} Every concept generically depends on some EventChain-record that is about it (Equation~\ref{eq:d6}); D6 is necessary for concept existence (it rules out ``disembodied'' concepts) and, together with D2, entails that a concept survives the destruction of any finite set of record copies provided at least one remains \cite{smith2012ontology}.

\textbf{Generic dependence of status on events.} A concept's status $s$ (e.g., \texttt{validated}) \emph{generically depends} on the events in its chain: $s$ depends on the existence of at least one \texttt{VALIDATE} event, but not on any \emph{specific} one (a different agent validating at a different time would still yield \texttt{validated}).

\textbf{Mutual dependence of relations on concepts.} A relation $r$ between concepts $c_1$ and $c_2$ \emph{mutually} (and rigidly) depends on both: $r$ cannot exist if either relatum ceases to exist. Formally, $\text{Exists}(r, t) \leftrightarrow \text{Exists}(c_1, t) \land \text{Exists}(c_2, t) \land \text{HoldsAt}(r, t)$, where $\text{Exists}(c, t)$ means $c$ has not been archived at $t$. If \texttt{concept-A} is deprecated, the relation ``\texttt{concept-A} isomorphic-to \texttt{concept-B}'' ceases to hold even if \texttt{concept-B} still exists---an ontological fact, not merely a query-side filter. The dependence is rigid because the relator cannot change its relata without becoming a different relator (Axiom~I7), consistent with UFO's relator theory \cite{ufo}.

We formalize these dependence patterns as an OWL~2 DL alignment fragment in the next subsection.

\subsection{Formal Alignment Fragment: OWL~2 DL Axiomatization}\label{subsec:formal-alignment}

We present the OWL~2 DL alignment fragment that formalizes the dependence patterns above in Appendix~\ref{app:owl-dl}. The revised fragment (v0.2) declares: (i)~core categories (\texttt{adl:Event} $\sqsubseteq$ \texttt{bfo:occurrent}, \texttt{adl:EventChain} $\sqsubseteq$ \texttt{bfo:process}, \texttt{adl:Concept} $\sqsubseteq$ \texttt{bfo:generically\_dependent\_continuant}); (ii)~L3 relation predicates as OWL object properties with symmetry/transitivity constraints (\texttt{isomorphic-to}, \texttt{specialisation-of}, \texttt{fork-of}, \texttt{related-to}, \texttt{analogical-to}, \texttt{co-occurs-with}, \texttt{mitigated-by}); (iii)~ontological dependence axioms (D2--D5) as object properties and disjointness constraints; and (iv)~illustrative SWRL rules for status-transition constraints. The fragment has been validated with ROBOT: OWL~2 DL profile conformance confirmed, structural consistency reported, and three SPARQL constraints verified on example documents (zero violations on valid data, correct detection on deliberately corrupted data). Full operational encoding of $\delta$ and $\gamma$ exceeds OWL~2 DL expressivity and is deferred to future work (FW1).

\textbf{SHACL/OWL expressivity boundaries.} The SHACL shape graph (Appendix~\ref{app:shacl}) validates L3 relation blocks against structural constraints: identifier format (\texttt{sh:pattern}), predicate membership (\texttt{sh:in}), confidence range (\texttt{sh:minInclusive}/\texttt{sh:maxInclusive}), and directionality (\texttt{sh:sparql}). Five of eight constraints use SHACL Core; three require SHACL-SPARQL. However, the status derivation function $\delta$ and the precondition checks are \emph{not expressible in SHACL or OWL~2 DL} for fundamental reasons: (i)~$\delta$ aggregates over event sequences (taking the maximum over lifecycle events), which requires recursion or fixed-point computation exceeding SHACL's declarative expressivity; (ii)~preconditions involve comparator dispatch over derived snapshots, which is procedural rather than declarative; (iii)~temporal reasoning (\texttt{HoldsAt}, \texttt{Revoked} conditions) requires time-indexed assertions over event sequences, which exceeds OWL~2 DL's snapshot-based semantics; (iv)~cryptographic integrity (SHA-256 correctness) is a computational property that cannot be expressed in description logic; (v)~\emph{cross-chain references} require accessing events outside the current chain, which violates the local-scope restriction of the precondition language (\S\ref{subsubsec:precondition-language}). Encoding cross-chain references in OWL~2~DL would require non-local scope during precondition evaluation, violating the referential transparency and deterministic evaluation guarantees of Theorem~8.

\textbf{Interoperability mapping table.} Supplementary Table S25 provides a normative mapping from ADL Lite core constructs to OWL~2 DL classes, SHACL shapes, and PROV-O/IAO properties. This mapping enables the applied ontology community to adopt ADL Lite's lifecycle semantics within existing RDF tooling without abandoning the runtime's derivation functions.

\textbf{Adoption path.} The bidirectional export/import (Turtle/RDF/XML via \texttt{owl\_export.py} / \texttt{jsonld\_export.py}) allows ADL Lite concepts to be consumed by RDF tooling without losing the EventChain structure. The exported RDF represents each event as a PROV-O activity with cryptographic linkage, and the concept as a PROV-O entity with derived status. This is a ``lossy but useful'' export: the lifecycle semantics ($\delta$, $\gamma$, preconditions) are not encoded in the RDF (they exceed OWL/SHACL expressivity), but the structural and provenance information is preserved. The SHACL shape graph (Appendix~\ref{app:shacl}) validates five of eight L3 structural constraints; the remaining three (directionality, cross-document existence, conditional symmetry) require SHACL-SPARQL. Full operational encoding of $\delta$ and preconditions in SHACL/OWL remains future work (FW6b).

\subsection{Comparison with Upper Ontologies: Alignment Choices, Deviations, and Costs}\label{subsec:alignment-deviations}

Table~\ref{tab:upper-ontology-mapping} identifies four intentional deviations from BFO/DOLCE/UFO, each with an interoperability cost. \textbf{Deviation 1 (no continuant--occurrent distinction at the language level):} BFO and DOLCE encode the distinction in the ontology language \cite{bfo,dolce}; ADL Lite enforces it conceptually through the event model, so its alignment is \emph{conceptual guidance} rather than strict commitment (no direct BFO module import or OWL-reasoner enforcement); OWL export therefore requires a translation layer. \textbf{Deviation 2 (no explicit quality hierarchy):} ADL Lite encodes qualities as payload fields (e.g., ``confidence'' carried by \texttt{VALIDATE} events), so DOLCE's quality taxonomy is not directly reusable without explicit mapping. \textbf{Deviation 3 (actions as first-class citizens):} actions are co-equal with objects at the language level, reflecting the event-first stance; direct reuse of UFO-B's process hierarchy is prevented, but UFO-B's ``Institutional Event'' category \cite{guizzardi2024ufob} remains compatible, and ADL Lite's precondition language is a \emph{variable-free ground fragment} (closed comparator enumeration, no quantification) that is a \emph{tractable specialisation} of UFO-B's FOL framework with $O(1)$ per-rule evaluation. \textbf{Deviation 4 (historical vs.\ ontological dependence):} the causal origin of a concept (traceable to the genesis event, an occurrent) is distinguished from its ontological bearer (the EventChain-record, an ICE): the concept is a GDC that depends on the record (D2), not on the process (D5), avoiding the category mistake of making a GDC existentially depend on an occurrent.

Supplementary Table S26 contrasts the operational mechanisms: ADL Lite embeds execution semantics \emph{inside} the data structure (the EventChain), whereas UFO-B delegates execution to an external workflow engine that interprets the ontology.

\textbf{Deviation 4: Historical vs. ontological dependence.} We distinguish the causal origin of a concept (traceable to the genesis event, an occurrent) from its ontological bearer (the EventChain-record, an ICE). The concept is a GDC that depends on the record (D2), not on the process (D5). This avoids the category mistake of making a GDC existentially depend on an occurrent. The ``historical'' or ``genetic'' relationship to the process is a causal narrative, not an ontological dependence in the Fine/BFO sense.
